% SPDX-License-Identifier: MIT
\chapter{Introduction}
\label{ch:introduction}

\app{} is a desktop application and Python library for visualising and annotating
phylogenetic trees. It runs entirely on your own machine: there is no server, no
account, no upload and no network access at any point while it runs.

\section{What it does}

\begin{description}[style=nextline,leftmargin=0pt,labelindent=0pt,labelwidth=0pt]
  \item[Reads what you already have.]
    Newick, New Hampshire eXtended (NHX), NEXUS and phyloXML. Parsing is
    deliberately forgiving --- real files break their specifications constantly,
    so \app{} loads them anyway and tells you what it had to interpret.
  \item[Draws it many ways.]
    Rectangular and slanted phylograms and cladograms, circular fans with a
    configurable arc and inner radius, radial layouts, and unrooted layouts by
    Felsenstein's equal-angle and equal-daylight algorithms.
  \item[Lets you edit the tree.]
    Reroot on any branch, by outgroup or at the midpoint; ladderize, rotate,
    collapse, prune and extract subtrees. Everything is undoable.
  \item[Annotates it.]
    Thirteen track types stack alongside the tree, described in
    Chapter~\ref{ch:tracks}.
  \item[Exports honestly.]
    SVG and PDF as vectors, PNG at any resolution. The on-screen canvas and the
    exported file are built from an identical description, so what you export is
    what you saw.
\end{description}

\section{Who it is for}

Anyone who has a tree and needs a figure of it: for a paper, a thesis, a talk or
a teaching handout. It is aimed particularly at people who cannot or would rather
not upload their data to a hosted service --- unpublished phylogenies, clinical
isolate data, or work under a data-transfer agreement. Nothing leaves your
machine, because the software has no way to send it anywhere.

The library half is aimed at a different need: regenerating a figure
reproducibly, from a script, as part of an analysis pipeline. Anything you can
build interactively you can also build in Python (Chapter~\ref{ch:python}) or
from the command line (Chapter~\ref{ch:cli}).

\section{Two packages}

\app{} ships as two Python packages, and the split is architectural rather than
legal --- both are MIT-licensed.

\begin{center}
\begin{tabular}{@{}lp{0.62\linewidth}@{}}
\toprule
\textbf{Package} & \textbf{Contents} \\
\midrule
\cmd{makeyourtree} &
  Parsers, tree model, layout engine, annotation model, scene graph, SVG
  renderer and the command-line interface. Imports no GUI toolkit, so it works
  headlessly in scripts, notebooks and continuous integration. Depends only on
  NumPy. \\
\cmd{makeyourtree\_studio} &
  The desktop application, built on PySide6/Qt~6. \\
\bottomrule
\end{tabular}
\end{center}

If you only want to generate figures from a script, you never need the second
package, and you can install without Qt entirely (Section~\ref{sec:install-library}).

\section{Three ideas behind the design}
\label{sec:design}

You do not need this section to use \app{}. It is here because it explains why
certain things behave the way they do.

\subsection{Two independent coordinates}

Every layout is built from two coordinates that are computed separately. The
\emph{cross} coordinate --- $y$ in a rectangular layout, angle in a circular one
--- comes only from the order of the tips. The \emph{along} coordinate --- $x$,
or radius --- comes only from branch length, or from depth in a cladogram.

The practical consequence is that the two things you do to a figure never
interfere. Ladderizing or rotating a clade changes the ordering and cannot alter
any distance; switching between a phylogram and a cladogram changes the distances
and cannot alter the ordering.

\subsection{Band space: why one track works in every layout}

Annotation tracks are not written for a particular layout. Each is written in
\emph{band space}, a coordinate system of \texttt{(row, offset)} where
\texttt{row} indexes the tip and \texttt{offset} is the distance outward from the
tree. A projector then maps band space into the page: a straight-line mapping for
linear layouts, a polar one for circular layouts.

This is why a heatmap becomes an annulus in a circular layout without you doing
anything, and why tangential labels stay the right way up on the left-hand side
of a fan. Figures~\ref{fig:combined-rectangular} and~\ref{fig:combined-circular}
are the same tracks with the same data; only the layout mode differs.

\subsection{One scene, many outputs}

Laying out and composing a figure produces a plain description of shapes and
text. The on-screen canvas draws that description, and so do the SVG, PDF and PNG
exporters. Because there is only one description, the screen and the file cannot
disagree.

\section{About this manual}

Chapters~\ref{ch:installation}--\ref{ch:quickstart} get you running.
Chapters~\ref{ch:interface}--\ref{ch:editing} cover the application.
Chapters~\ref{ch:annotation}--\ref{ch:tracks} cover annotation, and
Chapter~\ref{ch:tracks} is the reference for all thirteen track types.
Chapters~\ref{ch:export}--\ref{ch:python} cover output and automation.

Every figure in this manual was produced by \app{} from the synthetic example
data in the repository. Nothing is a mock-up.
